\subsection{Hilbert-Samuel series of a quotient module}

Lemma 5. — Let A be a ring, M an A-module, and \((\mathbf{P}_{n})\) , \((\mathbf{Q}_{n})\) two descending filtrations on M consisting of submodules. Suppose that we have \(P_{n} \supset Q_{n}\) and \(\text{long}_{\mathbf{A}}(\mathbf{P}_{n}/\mathbf{Q}_{n}) < +\infty\) for every \(n \in Z\) and that there exists an integer \(n_{1}\) such that \(Q_{n_{1}} = M\). In \(\mathbf{Z}((T))\), we have the inequalities

\[
\begin{array}{r l} \sum_ {n \in \mathbf {Z}} \mathrm{long} _ {\mathbf {A}} ((\mathrm{P} _ {n + 1} \cap \mathrm{Q} _ {n}) / \mathrm{Q} _ {n + 1}). & \mathrm{T} ^ {n} \leqslant \sum_ {n \in \mathbf {Z}} \mathrm{long} _ {\mathbf {A}} (\mathrm{P} _ {n + 1} / \mathrm{Q} _ {n + 1}). \\ & \mathrm{T} ^ {n} \leqslant \\ & \leqslant (1 - \mathrm{T}) ^ {- 1} \sum_ {n \in \mathbf {Z}} \mathrm{long} _ {\mathbf {A}} ((\mathrm{P} _ {n + 1} \cap \mathrm{Q} _ {n}) / \mathrm{Q} _ {n + 1}). \end{array}
\]

We need to prove that we have the inequalities

\begin{enumerate}
\def\labelenumi{(\arabic{enumi})}
\setcounter{enumi}{14}
\tightlist
\item
\end{enumerate}

\[
\operatorname{long} \left(\left(\mathrm{P} _ {n + 1} \cap \mathrm{Q} _ {n}\right) / \mathrm{Q} _ {n + 1}\right) \leqslant \operatorname{long} \left(\mathrm{P} _ {n + 1} / \mathrm{Q} _ {n + 1}\right),\tag{16}
\]

\[
\operatorname{long} \left(\mathrm{P} _ {n + 1} / \mathrm{Q} _ {n + 1}\right) \leqslant \sum_ {i \leqslant n} \operatorname{long} \left(\left(\mathrm{P} _ {i + 1} \cap \mathrm{Q} _ {i}\right) / \mathrm{Q} _ {i + 1}\right).
\]

The first is evident. On the other hand, we have \(P_{n+1} \cap Q_i = P_{n+1}\) for \(i \leqslant n_1\) and \(P_{n+1} \cap Q_{n+1} = Q_{n+1}\) ; we deduce the inequality

\[
\operatorname{long} \left(\mathrm{P} _ {n + 1} / \mathrm{Q} _ {n + 1}\right) \leqslant \sum_ {i \leqslant n} \operatorname{long} \left(\left(\mathrm{P} _ {n + 1} \cap \mathrm{Q} _ {i}\right) / \left(\mathrm{P} _ {n + 1} \cap \mathrm{Q} _ {i + 1}\right)\right).
\]

But the A-module \((\mathrm{P}_{n+1} \cap \mathrm{Q}_i)/(\mathrm{P}_{n+1} \cap \mathrm{Q}_{i+1})\) is isomorphic to a submodule of \((\mathrm{P}_{i+1} \cap \mathrm{Q}_i)/\mathrm{Q}_{i+1}\) , and inequality (16) follows.

Lemma 6. --- Let A be a ring, M an A-module, and \((F_{n})\) a decreasing filtration on M consisting of submodules; suppose there exists an integer \(n_{1}\) such that \(F_{n_{1}} = M\). Let f be an endomorphism of M, \(M'\) its kernel, and \(M''\) its cokernel. We endow \(M'\) with the filtration \((F_{n}')\) induced by \((F_{n})\) and \(M''\) with the quotient filtration \((F_{n}'')\) of \((F_{n})\). We assume that \(F_{n}/F_{n+1}\) is of finite length for every \(n \in \mathbf{Z}\) and that there exists an integer \(\delta\) such that \(f(F_{n}) \subset F_{n+\delta}\). Let \(\varphi\) be the graded endomorphism of degree \(\delta\) of the graded module \(\operatorname{gr}(M) = \bigoplus_{n} F_{n}/F_{n+1}\) induced by f. Between the following elements of \(\mathbf{Z}((T))\)

\[
\mathrm{H} _ {\mathbf {M}} = \sum_ {n \in \mathbf {Z}} \operatorname{long} _ {\mathbf {A}} (\mathrm{F} _ {n} / \mathrm{F} _ {n + 1}). \mathrm{T} ^ {n}
\]

\[
\mathrm{H} _ {\mathbf {M} ^ {\prime}} = \sum_ {n \in \mathbf {Z}} \operatorname{long} _ {\mathbf {A}} (\mathrm{F} _ {n} ^ {\prime} / \mathrm{F} _ {n + 1} ^ {\prime}). \mathrm{T} ^ {n}
\]

\[
\mathrm{H} _ {\mathrm{M} ^ {\prime \prime}} = \sum_ {n \in \mathbf {Z}} \operatorname{long} _ {\mathrm{A}} (\mathrm{F} _ {n} ^ {\prime \prime} / \mathrm{F} _ {n + 1} ^ {\prime \prime}). \mathrm{T} ^ {n}
\]

\[
\mathrm{P} _ {\operatorname{Ker} (\varphi)} = \sum_ {n \in \mathbf {Z}} \operatorname{long} _ {\mathrm{A}} (\operatorname{Ker} (\varphi_ {n})). \mathrm{T} ^ {n},
\]

we have the inequalities

\begin{enumerate}
\def\labelenumi{(\arabic{enumi})}
\setcounter{enumi}{16}
\tightlist
\item
\end{enumerate}

\[
\mathrm{H} _ {\mathbf {M} ^ {\prime}} \leqslant \mathrm{P} _ {\mathbf {K e r} (\varphi)}\tag{18}
\]

\[
(1 - \mathrm{T} ^ {\delta}). \mathrm{H} _ {\mathbf {M}} ^ {(1)} + \mathrm{T} ^ {\delta}. \mathrm{P} _ {\operatorname{Ker} (\varphi)} \leqslant \mathrm{H} _ {\mathbf {M} ^ {\prime \prime}} ^ {(1)} \leqslant (1 - \mathrm{T} ^ {\delta}). \mathrm{H} _ {\mathbf {M}} ^ {(1)} + \mathrm{T} ^ {\delta}. \mathrm{P} _ {\operatorname{Ker} (\varphi)} ^ {(1)}.
\]

The sequence of submodules \(\mathbf{G}_{n}=f^{-1}(\mathbf{F}_{n+\delta})\) on M is a decreasing filtration, and we have \(F_{n}\subset G_{n}\) for every integer n.

By definition, we have \(\operatorname{Ker}(\varphi_n) = (\mathbf{G}_{n+1} \cap \mathbf{F}_n)/\mathbf{F}_{n+1}\) , whence

\[
\mathrm{P} _ {\operatorname{Ker} (\varphi)} = \sum_ {n \in \mathbf {Z}} \operatorname{long} _ {\mathrm{A}} ((\mathrm{G} _ {n + 1} \cap \mathrm{F} _ {n}) / \mathrm{F} _ {n + 1}). \mathrm{T} ^ {n}.\tag{19}
\]

For every \(n\) , the A-module \((\mathbf{M}' \cap \mathbf{F}_n) / (\mathbf{M}' \cap \mathbf{F}_{n+1})\) is identified with a submodule of \((\mathbf{G}_{n+1} \cap \mathbf{F}_n) / \mathbf{F}_{n+1}\) , and inequality (17) follows immediately from (19). Moreover, by Lemma 5, we have

\[
\mathrm{P} _ {\operatorname{Ker} (\varphi)} \leqslant \sum_ {n \in \mathbf {Z}} \operatorname{long} _ {\mathrm{A}} \left(\mathrm{G} _ {n + 1} / \mathrm{F} _ {n + 1}\right). \mathrm{T} ^ {n} \leqslant \mathrm{P} _ {\operatorname{Ker} (\varphi)} ^ {(1)}.\tag{20}
\]

For every \(n \in \mathbf{Z}\) we have an exact sequence of A-modules

\[
0 \longrightarrow \mathrm{G} _ {n + 1} / \mathrm{F} _ {n + 1} \longrightarrow \mathrm{M} / \mathrm{F} _ {n + 1} \xrightarrow {f _ {n}} \mathrm{M} / \mathrm{F} _ {n + \delta + 1} \longrightarrow \mathrm{M} ^ {\prime \prime} / \mathrm{F} _ {n + \delta + 1} ^ {\prime \prime} \longrightarrow 0,
\]

where \(f_{n}\) is deduced from \(f\) by passing to quotients. Consequently, we have

\[
\operatorname{long} _ {\mathrm{A}} \left(\mathrm{M} ^ {\prime \prime} / \mathrm{F} _ {n + \delta + 1} ^ {\prime \prime}\right) = \operatorname{long} _ {\mathrm{A}} \left(\mathrm{M} / \mathrm{F} _ {n + \delta + 1}\right) - \operatorname{long} _ {\mathrm{A}} \left(\mathrm{M} / \mathrm{F} _ {n + 1}\right) + \operatorname{long} _ {\mathrm{A}} \left(\mathrm{G} _ {n + 1} / \mathrm{F} _ {n + 1}\right).
\]

Multiplying by \(T^{n+\delta}\) and summing over n, we obtain

\[
\mathrm{H} _ {\mathbf {M} ^ {\prime \prime}} ^ {(1)} = (1 - \mathrm{T} ^ {\delta}) \mathrm{H} _ {\mathbf {M}} ^ {(1)} + \mathrm{T} ^ {\delta}. \sum_ {n \in \mathbf {Z}} \operatorname{long} _ {\mathrm{A}} (\mathrm{G} _ {n + 1} / \mathrm{F} _ {n + 1}). \mathrm{T} ^ {n},\tag{21}
\]

and inequality (18) follows immediately from (20) and (21).

Lemma 7. — We keep the notations of Lemma 6.

\begin{enumerate}
\def\labelenumi{\alph{enumi})}
\item
  We have the inequality \(\mathbf{H}_{\mathbf{M}^{\prime \prime}}^{(1)}\geqslant \frac{1 - \mathrm{T}^{\delta}}{1 - \mathrm{T}}\mathbf{H}_{\mathbf{M}}.\)
\item
  For equality to hold, it is necessary and sufficient that φ be injective.
\item
  If this is the case, we have \(\mathbf{M}^{\prime}\subset \bigcap_{n}\mathbf{F}_{n}\) and the sequence of A-modules
\end{enumerate}

\[
0 \longrightarrow \operatorname{gr} (\mathrm{M}) \stackrel {\varphi} {\longrightarrow} \operatorname{gr} (\mathrm{M}) \stackrel {v} {\longrightarrow} \operatorname{gr} (\mathrm{M} ^ {\prime \prime}) \longrightarrow 0,
\]

where v is the canonical map, is exact.

Assertions a) and b) follow from formula (18) of Lemma 6 and from the definition \(\mathbf{H}_{\mathbf{M}}^{(1)} = (1 - \mathbf{T})^{-1}.\mathbf{H}_{\mathbf{M}}\) .

Suppose that \(\varphi\) is injective. By III, § 2, no. 8, Th. 1, (i), we have

\[
\operatorname{Ker} (f) \subset f ^ {- 1} (\mathrm{F} _ {n + \delta}) = \mathrm{F} _ {n}
\]

for all n, whence the first assertion of c). Moreover, we have an exact sequence

\[
0 \longrightarrow \mathrm{M} / \mathrm{M} ^ {\prime} \stackrel {f ^ {\prime}} {\longrightarrow} \mathrm{M} \longrightarrow \mathrm{M} / f (\mathrm{M}) \longrightarrow 0,
\]

where \(f'\) is deduced from \(f\) by passing to the quotient. If \(\varphi\) is injective, we have as above \(f'^{-1}(\mathbf{F}_n) = \mathbf{F}_{n-\delta}/\mathbf{M}'\). Thus the filtration on \(\mathbf{M}/\mathbf{M}'\) deduced by inverse image under \(f'\) of the filtration \(\mathbf{F}\) on \(\mathbf{M}\) is the filtration \(n \mapsto \mathbf{F}_{n-\delta}/\mathbf{M}'\); the associated graded is \(\operatorname{gr}(\mathbf{M})(-\delta)\) and we have an exact sequence of graded modules (III, § 2, no. 4, prop. 2)

\[
0 \longrightarrow \operatorname{gr} (\mathrm{M}) (- \delta) \xrightarrow {\varphi^ {\prime}} \operatorname{gr} (\mathrm{M}) \longrightarrow \operatorname{gr} \left(\mathrm{M} ^ {\prime \prime}\right) \longrightarrow 0,
\]

where \(\varphi_{n}^{\prime} = \varphi_{n - \delta}\) for every \(n\). This completes the proof of \(c\) ).

PROPOSITION 6. — Let A be a Noetherian ring, M a finitely generated A-module, and q an ideal of A such that M/qM has finite length. Let F be a q-good filtration of M, and let gr(A) = \(\bigoplus_{n\geqslant 0}(\mathfrak{q}^n /\mathfrak{q}^{n + 1})\) the graded ring associated with A for the q-adic filtration.

Let \((x_{1},\ldots,x_{s})\) be a sequence of elements of A, \((\delta_{1},\ldots,\delta_{s})\) a sequence of strictly positive integers such that \(x_{i}\in\mathfrak{q}^{\delta_{i}}\) for \(1\leqslant i\leqslant s\) , and let \(\xi_{i}\) be the class of \(x_{i}\) in \(\mathrm{gr}_{\delta_{i}}(\mathbf{A})=\mathfrak{q}^{\delta_{i}}/\mathfrak{q}^{\delta_{i}+1}\) .

\begin{enumerate}
\def\labelenumi{\alph{enumi})}
\tightlist
\item
  Endow the A-module \(\overline{\mathbf{M}} = \mathbf{M}/(x_1\mathbf{M} + \cdots + x_s\mathbf{M})\) with the quotient q-good filtration \(\overline{\mathbf{F}}\) of F. We then have in \(\mathbf{Z}((\mathbf{T}))\) the inequality
\end{enumerate}

\[
\mathrm{H} _ {\mathrm{M}, \overline {{\mathrm{F}}}} ^ {(s)} \geqslant \left(\prod_ {i = 1} ^ {s} \frac {1 - \mathrm{T} ^ {\delta_ {i}}}{1 - \mathrm{T}}\right). \mathrm{H} _ {\mathrm{M}, \mathrm{F}}.\tag{22}
\]

\begin{enumerate}
\def\labelenumi{\alph{enumi})}
\setcounter{enumi}{1}
\tightlist
\item
  For equality to hold in (22), it is necessary and sufficient that the sequence \((\xi_1, \ldots, \xi_s)\) of elements of the ring gr(A) be completely secant for the module gr(M) = \(\bigoplus_{n} (\mathbf{F}_n / \mathbf{F}_{n+1})\) .
\end{enumerate}

In this case, the canonical homomorphism from \(\mathrm{gr}(\mathbf{M}) / \sum_{i=1}^{s} \xi_i \cdot \mathrm{gr}(\mathbf{M})\) into \(\mathrm{gr}(\overline{\mathbf{M}}) = \bigoplus_{n} (\overline{\mathbf{F}}_n / \overline{\mathbf{F}}_{n+1})\) is an isomorphism.

\begin{enumerate}
\def\labelenumi{\alph{enumi})}
\setcounter{enumi}{2}
\tightlist
\item
  Suppose the conditions of b) are satisfied, and that each of the A-modules \(\mathbf{M}_{i} = \mathbf{M}/(x_{1}\mathbf{M} + \cdots + x_{i}\mathbf{M})\) for \(0 \leqslant i < s\) is separated for the q-adic topology \footnote{This occurs in particular if \(q\) is contained in the radical of \(A\) (III, § 3, no. 3, Prop. 6).}. Then the sequence \((x_{1}, ..., x_{s})\) is completely secant for the A-module M.
\end{enumerate}

When \(s = 1\) , one has \(\bigcap_{n} F_{n} = \bigcap_{n} q^{n} M\) and the sequence \(\{\xi_{1}\}\) is completely secant for \(\operatorname{gr}(M)\) if and only if the homothety with ratio \(\xi_{1}\) in \(\operatorname{gr}(M)\) is injective. Prop. 6 then follows immediately from Lemma 7 applied to the homothety \(f = (x_{1})_{M}\) in \(M\) .

Suppose that we have \(s \geqslant 2\) and reason by induction on \(s\). The induction hypothesis applied to the A-module \(\mathbf{M}_{1} = \mathbf{M}/x_{1}\mathbf{M}\) endowed with the quotient filtration G of F, and to the sequence \((x_{2}, ..., x_{s})\) yields the inequality

\[
\mathrm{H} _ {\overline {{\mathrm{M,F}}}} ^ {(s - 1)} \geqslant \left(\prod_ {i = 2} ^ {s} \frac {1 - \mathrm{T} ^ {\delta_ {i}}}{1 - \mathrm{T}}\right). \mathrm{H} _ {\mathrm{M} _ {1}, \mathrm{G}};\tag{23}
\]

there is equality if and only if the sequence \((\xi_{2}, \ldots, \xi_{s})\) is completely secant for the \(\operatorname{gr}(A)\)-module \(\operatorname{gr}(\mathbf{M}_{1}) = \bigoplus_{n} \mathrm{G}_{n}/\mathrm{G}_{n+1}\). Since the elements \(\frac{1 - \mathrm{T}^{\delta_{i}}}{1 - \mathrm{T}}\) of \(\mathbf{Z}((\mathrm{T}))\) are positive, the case \(s = 1\) already treated and formula (23) yield the inequalities

\[
\mathrm{H} _ {\mathrm{M}, \mathrm{F}} ^ {(s)} \geqslant \left(\prod_ {i = 2} ^ {s} \frac {1 - \mathrm{T} ^ {\delta_ {i}}}{1 - \mathrm{T}}\right). \mathrm{H} _ {\mathrm{M} _ {1}, \mathrm{G}} ^ {(1)} \geqslant \left(\prod_ {i = 1} ^ {s} \frac {1 - \mathrm{T} ^ {\delta_ {i}}}{1 - \mathrm{T}}\right). \mathrm{H} _ {\mathrm{M}, \mathrm{F}}.\tag{24}
\]

This proves a).

One can have equality in (22) only if one has simultaneously equality in (23) and equality

\[
\mathrm{H} _ {\mathbf {M} _ {1}, \mathbf {G}} ^ {(1)} = \left(\frac {1 - \mathrm{T} ^ {\delta_ {1}}}{1 - \mathrm{T}}\right) \cdot \mathrm{H} _ {\mathbf {M}, \mathbf {F}}.\tag{25}
\]

This latter relation means that \(\{\xi_{1}\}\) is completely secant for \(\operatorname{gr}(M)\) and implies that the canonical homomorphism from \(\operatorname{gr}(M)/\xi_{1}\operatorname{gr}(M)\) into \(\operatorname{gr}(M_{1})\) is an isomorphism. In other words, equality holds in (22) if and only if \(\{\xi_{1}\}\) is completely secant for \(\operatorname{gr}(M)\) and \(\{\xi_{2}, ..., \xi_{s}\}\) is completely secant for \(\operatorname{gr}(M)/\xi_{1}\operatorname{gr}(M)\). This means that \(\{\xi_{1}, ..., \xi_{s}\}\) is completely secant for \(\operatorname{gr}(M)\) (A, X, p. 160, cor. 2). We have thus proved the equivalence of the two conditions of \(b\)). Assume they are satisfied; then, by the induction hypothesis, \(\operatorname{gr}(M)\) is identified with \(\operatorname{gr}(M_{1})/\sum_{i=2}^{s} \xi_{i}\operatorname{gr}(M_{1})\); moreover, since \(\operatorname{gr}(M_{1})\) is identified with \(\operatorname{gr}(M)/\xi_{1}\operatorname{gr}(M)\), the last assertion of \(b\)) is thus satisfied.

Suppose now that \(\{\xi_1, \ldots, \xi_s\}\) is completely secant for \(\operatorname{gr}(\mathbf{M})\) and that each \(\mathbf{M}_i\) is separated for the q-adic topology (for \(0 \leqslant i < s\) ). From the above and the induction hypothesis, the sequence \((x_2, \ldots, x_s)\) is completely secant for \(\mathbf{M}_1\); since we have \(\mathbf{M}_1 = \mathbf{M}/x_1\mathbf{M}\) and \(\{x_1\}\) is completely secant for \(\mathbf{M}\), the sequence \((x_1, x_2, \ldots, x_s)\) is completely secant for \(\mathbf{M}\) (A, X, p. 160, th. 1).

